\chapter{Core Invariants}

\section{Role of the Core Invariants}

The Unified Rigidity Framework uses a small family of invariants to compare
different systems without pretending that the systems are identical. A graph
refinement process, a spectral operator, a physical energy functional, and a
capacity-bounded computation may have different native languages, but each can
be audited by asking the same structural questions:

\begin{enumerate}
\item What is the configuration space?
\item What uncertainty or defect remains unresolved?
\item How much information can one admissible step remove?
\item What counts as terminal or defect-free?
\item What coercive quantity prevents free collapse?
\item What lower bound on refinement depth follows?
\end{enumerate}

This chapter fixes the basic invariants used in later chapters. The definitions
are intentionally abstract, because the framework is meant to support multiple
instantiations. A later theorem must always specify the concrete instance being
used.

\section{Capacity--Entropy Relation}

Let \(X\) be a state space and let \(H:X\to\mathbb R_{\ge 0}\) be an entropy or
defect-size functional. Let \(C>0\) denote the maximum amount of relevant
information that one admissible refinement step can remove.

\begin{definition}[One-step capacity gap]
The one-step capacity gap is
\[
\Phi_C(x)=H(x)-C.
\]
It measures how much unresolved information remains beyond what one admissible
step can remove.
\end{definition}

The inequality \(\Phi_C(x)\le 0\) should not be read as a universal
realizability condition for all states. It is a one-step collapse condition.
If \(H(x)>C\), then the state may still be admissible, but it cannot be fully
resolved in one capacity-bounded step.

\begin{theorem}[Capacity--Entropy Step Bound]
Suppose each admissible refinement step reduces entropy by at most \(C\). If
\(x_0,x_1,\ldots,x_t\) is a refinement trajectory, then
\[
H(x_0)-H(x_t)\le tC.
\]
\end{theorem}

\begin{proof}
For each step, the capacity hypothesis gives
\[
H(x_i)-H(x_{i+1})\le C.
\]
Summing this inequality from \(i=0\) to \(t-1\) gives
\[
\sum_{i=0}^{t-1}\bigl(H(x_i)-H(x_{i+1})\bigr)\le tC.
\]
The left side telescopes to \(H(x_0)-H(x_t)\), proving the claim.
\end{proof}

\begin{theorem}[Depth lower bound]
If the terminal condition requires \(H(x_t)=0\), then every admissible
trajectory from \(x_0\) to a terminal state satisfies
\[
t\ge \frac{H(x_0)}{C}.
\]
\end{theorem}

\begin{proof}
Set \(H(x_t)=0\) in the preceding theorem. Then \(H(x_0)\le tC\), so
\(t\ge H(x_0)/C\).
\end{proof}

\section{Entropy}

Entropy is the invariant that records unresolved multiplicity.

\begin{definition}[Entropy functional]
An entropy functional on a state space \(X\) is a map
\[
H:X\to\mathbb R_{\ge 0}
\]
such that \(H(x)=0\) represents a terminal, resolved, or defect-free state,
relative to the chosen model.
\end{definition}

In a finite configuration model, if \(x\) represents a set \(\Omega_x\) of
still-possible configurations, the basic entropy is
\[
H(x)=\log |\Omega_x|.
\]
In a probabilistic model with distribution \(p\), the standard entropy is
\[
H(p)=-\sum_y p(y)\log p(y).
\]

The framework does not require every domain to use the same entropy formula.
It requires that the entropy be explicitly declared and that its behavior under
refinement be proved or assumed as a named boundary.

\begin{example}[Finite search space]
If a system begins with \(2^n\) possible configurations, then
\[
H_0=\log_2(2^n)=n.
\]
If each admissible step removes at most one bit of relevant uncertainty, then
at least \(n\) steps are required to isolate a single configuration.
\end{example}

\section{Capacity}

Capacity is the invariant that records the information throughput of an
admissible step.

\begin{definition}[Step capacity]
A refinement class \(\mathcal R\) has step capacity \(C\) with respect to
entropy \(H\) if every admissible refinement step \(R\in\mathcal R\) satisfies
\[
H(x)-H(Rx)\le C.
\]
\end{definition}

Capacity is relative to the admissible class. If the class allows a global
oracle, then \(C\) may be large. If the class only allows bounded local
inspection, then \(C\) is controlled by the local view. This distinction is
essential: URF lower bounds apply only to the refinement class for which the
capacity estimate has been established.

\begin{example}[Local graph update]
For a bounded-degree graph with fixed radius \(R\), a vertex update sees only
a bounded neighborhood type. If the degree bound and radius are fixed, then
the number of possible local views is bounded independently of the total graph
size. This gives a candidate source of bounded per-step information.
\end{example}

\section{Depth}

Depth measures the number of admissible steps required to reach a terminal
state.

\begin{definition}[EntropyDepth]
Let \(T\subseteq X\) be the terminal set. For an admissible refinement process
\(x_{t+1}=R_t(x_t)\), define
\[
\ED(x_0)=\inf\{t\ge 0:x_t\in T\}.
\]
When \(T=\{x:H(x)=0\}\), this is the entropy-collapse depth.
\end{definition}

Depth is not merely runtime. It is structural length inside a specified
refinement model. A fast algorithm outside the admissible model does not
contradict a depth lower bound inside the model.

\begin{theorem}[Linear-depth criterion]
Suppose \(H(x_n)\ge an\) for a family of inputs \(x_n\), with \(a>0\), and
suppose every admissible step removes at most \(C\) entropy, with \(C\)
independent of \(n\). Then
\[
\ED(x_n)=\Omega(n).
\]
\end{theorem}

\begin{proof}
By the depth lower bound,
\[
\ED(x_n)\ge \frac{H(x_n)}{C}\ge \frac{a}{C}n.
\]
Thus the depth grows at least linearly in \(n\).
\end{proof}

\section{Spectral Gap}

Spectral gap is the analytic form of rigidity used when the state space is
linear or operator-theoretic.

\begin{definition}[Spectral gap]
Let \(L\) be a nonnegative self-adjoint operator. The spectral gap above the
kernel is the largest \(\lambda_1\ge 0\) such that
\[
\langle f,Lf\rangle\ge \lambda_1\|f\|^2
\]
for all \(f\perp\ker L\).
\end{definition}

A positive spectral gap prevents low-cost deformation away from the kernel.
In URF language, it supplies a coercive mechanism: unresolved defect cannot
vanish without paying a controlled amount.

\begin{example}[Cycle graph]
For the cycle graph \(C_n\), the graph Laplacian eigenvalues are
\[
\lambda_k=2-2\cos\left(\frac{2\pi k}{n}\right).
\]
The first nonzero eigenvalue is
\[
\lambda_1=2-2\cos\left(\frac{2\pi}{n}\right).
\]
As \(n\) grows, this gap tends to zero, so cycle graphs do not form a uniform
spectral expander family.
\end{example}

\section{Invariant Checklist}

A URF argument should identify the following data:

\begin{center}
\begin{tabular}{ll}
\toprule
Object & Required role \\
\midrule
\(X\) & admissible configuration space \\
\(\mathcal R\) & admissible refinement class \\
\(H\) & entropy or unresolved defect \\
\(C\) & per-step capacity bound \\
\(T\) & terminal set \\
\(F\) & rigidity or coercivity functional \\
\(\lambda_1\) & spectral gap when an operator model is present \\
\bottomrule
\end{tabular}
\end{center}

If any entry is missing, the argument is not yet a theorem-level URF
instantiation. It may still be a useful target, roadmap, or frontier surface.

\begin{remark}[Boundary]
This chapter defines the invariant vocabulary and proves only the elementary
capacity-depth bookkeeping statements above. It does not prove unrestricted
Chronos-RR, unrestricted H4.1/FGL, P vs NP, or any Clay problem.
\end{remark}
